% Part: first-order-logic
% Chapter: tableaux
% Section: quantifier-rules

\documentclass[../../../include/open-logic-section]{subfiles}

\begin{document}

\olfileid[fr]{fol}{tab}{qrl}

\olsection{Règles pour les quantificateurs}

\subsection{Règles pour $\lforall$}

\begin{defish}
\AxiomC{\sFmla{\True}{\lforall[x][!A(x)]}}
\RightLabel{\TRule{\True}{\forall}}
\UnaryInfC{\sFmla{\True}{!A(t)}}
\DisplayProof
\hfill
\AxiomC{\sFmla{\False}{\lforall[x][!A(x)]}}
\RightLabel{\TRule{\False}{\lforall}}
\UnaryInfC{\sFmla{\False}{!A(a)}}
\DisplayProof
\end{defish}

Dans \TRule{\True}{\lforall}, $t$ est un terme clos, c'est-à-dire
sans variables. Dans \TRule{\False}{\lforall}, $a$ est
!!a{constant} qui ne doit figurer nulle part dans la branche au-dessus
de l'application de la règle. On appelle $a$ l'\emph{eigenvariable}
de l'inférence \TRule{\False}{\forall}.\footnote{Nous employons le
terme « eigenvariable » bien que $a$ soit !!a{constant} dans la
règle ci-dessus, pour des raisons historiques.}

\subsection{Règles pour $\lexists$}

\begin{defish}
\AxiomC{\sFmla{\True}{\lexists[x][!A(x)]}}
\RightLabel{\TRule{\True}{\lexists}}
\UnaryInfC{\sFmla{\True}{!A(a)}}
\DisplayProof
\hfill
\AxiomC{\sFmla{\False}{\lexists[x][!A(x)]}}
\RightLabel{\TRule{\False}{\lexists}}
\UnaryInfC{\sFmla{\False}{!A(t)}}
\DisplayProof
\end{defish}

Ici encore, $t$ est un terme clos, et $a$ est !!a{constant} qui ne
figure pas dans la branche au-dessus de l'application de la règle
\TRule{\True}{\lexists}. On appelle $a$ l'\emph{eigenvariable} de
l'inférence \TRule{\True}{\lexists}.

La condition selon laquelle une eigenvariable ne doit pas figurer
dans la branche au-dessus d'une inférence \TRule{\False}{\lforall}
ou \TRule{\True}{\lexists} s'appelle la \emph{condition sur
l'eigenvariable}.

\begin{explain}
Rappelons notre convention : lorsque $!A$ est !!a{formula} dans
laquelle la !!{variable}~$x$ est libre, nous l'indiquons en écrivant
$!A(x)$. Dans ce contexte, $!A(t)$ abrège alors~$\Subst{!A}{t}{x}$.
On pourrait donc aussi écrire la règle \TRule{\False}{\lexists}
comme suit :
\begin{prooftree}
  \AxiomC{\sFmla{\False}{\lexists[x][!A]}}
  \RightLabel{\TRule{\False}{\lexists}}
  \UnaryInfC{\sFmla{\False}{\Subst{!A}{t}{x}}}
\end{prooftree}
Le terme $t$ peut déjà figurer dans~$!A$ ; par exemple, $!A$ peut
être~$\Atom{\Obj P}{t,x}$. Inférer
$\sFmla{\False}{\Atom{\Obj P}{t,t}}$ de
$\sFmla{\False}{\lexists[x][\Atom{\Obj P}{t,x}]}$ constitue donc
une application correcte de~\TRule{\False}{\lexists}. En revanche,
les conditions sur l'eigenvariable dans \TRule{\False}{\lforall}
et~\TRule{\True}{\lexists} imposent notamment que la
!!{constant}~$a$ ne figure pas dans~$!A$. On ne peut donc pas inférer
correctement $\sFmla{\False}{\Atom{\Obj P}{a,a}}$ de
$\sFmla{\False}{\lforall[x][\Atom{\Obj P}{a,x}]}$ au moyen
de~$\TRule{\False}{\lforall}$.
\end{explain}

\begin{explain}
Dans \TRule{\True}{\lforall} et \TRule{\False}{\lexists}, le terme
$t$ n'est soumis à aucune restriction supplémentaire : il doit
seulement être clos. En revanche, dans les règles
\TRule{\True}{\lexists} et \TRule{\False}{\lforall}, la condition
sur l'eigenvariable exige que la !!{constant}~$a$ ne figure nulle part
dans la branche au-dessus de l'inférence correspondante. Cette
condition est nécessaire à la correction du système. Sans elle, le
tableau suivant serait un !!{tableau} fermé pour
$\lexists[x][\formula{A}(x)] \lif \lforall[x][\formula{A}(x)]$ :
\begin{center}
\begin{tableau}{}
  [\sFmla{\False}{\lexists[x][\formula{A}(x)] \lif \lforall[x][\formula{A}(x)]}, just=\TAss
    [\sFmla{\True}{\lexists[x][\formula{A}(x)]},
      just={\TRule{\False}{\lif}[1]}
      [\sFmla{\False}{\lforall[x][\formula{A}(x)]},
        just={\TRule{\False}{\lif}[1]}
        [\sFmla{\True}{\formula{A}(a)},
          just={\TRule{\True}{\lexists}[2]}
          [\sFmla{\False}{\formula{A}(a)},
            just={\TRule{\False}{\lforall}[3]}, close]
        ]
      ]
    ]
  ]
\end{tableau}
\end{center}
Or $\lexists[x][\formula{A}(x)] \lif
\lforall[x][\formula{A}(x)]$ n'est pas valide.
\end{explain}

\begin{editorial}
Note de l'édition française : l'original exige d'abord que $t$ soit
clos, puis affirme qu'il n'est soumis à aucune restriction. La seconde
occurrence est précisée ici : il n'existe aucune restriction
supplémentaire à l'exigence de fermeture.
\end{editorial}

\end{document}
